\section{Finite jet tests and initial monomials}\label{sec:jets}

We first explain how an injective jet test for initial monomials gives an
injective jet test for the original holomorphic functions.  The argument
uses finite matrices and convergent power series.

For a holomorphic function $f$ near a point $q$ of a smooth complex surface,
write $\ord_q(f)$ for its order of vanishing, with $\ord_q(0)=\infty$.
In local coordinates, $\ord_q(f)\ge m$ means that all partial derivatives of
total order less than $m$ vanish at $q$.  For distinct points
$\boldsymbol q=(q_1,\ldots,q_r)$ and a finite-dimensional space $F$ of
functions defined near them, its \emph{order-$m$ jet map} is
\[
 J^m_{\boldsymbol q}\colon F\longrightarrow
 \bigoplus_{j=1}^r\mathcal O_{q_j}/\mathfrak m_{q_j}^{m},
 \qquad f\longmapsto (f\bmod\mathfrak m_{q_j}^{m})_{j=1}^r.
\]
Here $\mathcal O_{q_j}$ is the ring of holomorphic germs at $q_j$, and
$\mathfrak m_{q_j}$ is its ideal of germs vanishing there.  We call the
corresponding jet test \emph{injective} when this map is injective, or
equivalently when no nonzero $f\in F$ vanishes to order at least $m$ at
every $q_j$.  In two local coordinates, the matrix of this map records
the derivatives of total order less than $m$ of a basis of $F$; it has
$r\binom{m+1}{2}$ rows.  Its injectivity is exactly full column rank.

The same definition applies to sections of a line bundle using local
frames.  Multiplication by a holomorphic unit induces an invertible map
on each jet quotient, so the definition is independent of those frames.
A biholomorphic change of coordinates preserves the powers of the
maximal ideal and hence preserves orders of vanishing and jet
injectivity.

Finite matrices also underlie monomial diagram methods for plane
interpolation \cite[Propositions~10 and~13]{Dumnicki2007}. The transfer
below applies to convergent local functions and follows from persistence
of one nonzero maximal minor.

\begin{lemma}[Persistence of a finite jet test]\label{lem:rank-persistence}
Let $m,r,N$ be positive integers and let $q_1,\ldots,q_r$ be distinct
points of an open subset $U\subset\C^2$.  For $1\le i\le N$ and
$0\le s<s_0$, let $f_{i,s}$ be holomorphic on $U$.  Suppose that, as
$s\downarrow0$, every derivative of $f_{i,s}$ of total order less than
$m$ at each $q_j$ converges to the corresponding derivative of $f_{i,0}$.
If the jet matrix of $f_{1,0},\ldots,f_{N,0}$ has column rank $N$, then
the jet matrix of $f_{1,s},\ldots,f_{N,s}$ also has column rank $N$ for
all sufficiently small positive $s$.

In particular, convergence locally uniformly on $U$ suffices for the
derivative hypothesis.
\end{lemma}

\begin{proof}
A nonzero $N\times N$ minor of the limiting jet matrix remains nonzero
for all sufficiently small $s$, since each of its entries converges.
The final assertion follows from the Cauchy integral formula on small
polydiscs around the finitely many points.
\end{proof}

For a bounded set $\Omega\subset\R^2$, denote by
\[
 \mathcal M(\Omega)
 =\Span_{\C}\{x^\alpha z^\beta:
                 (\alpha,\beta)\in\Omega\cap\Z^2\}
\]
the finite-dimensional space of Laurent polynomials on the torus
$\T$.  When another pair of torus coordinates is specified,
the same notation uses that pair in place of $(x,z)$.

We now consider convergent power series at the origin.  Fix $t>0$ and
order pairs $(\alpha,\beta)\in\Z_{\ge0}^2$ lexicographically by
\[
 (\alpha+t\beta,\beta).
\]
Every bounded interval of weights contains only finitely many exponent
pairs.  Thus each nonzero series $f\in\C\{x,z\}$ has a first nonzero
exponent in this order, denoted by $\nu_t(f)$.  For a finite-dimensional
subspace $F\subset\C\{x,z\}$, put
\[
 E_t(F)=\{\nu_t(f):0\ne f\in F\}.
\]

\begin{proposition}[Transfer from initial monomials]
\label{prop:initial-transfer}
Let $F\subset\C\{x,z\}$ be a finite-dimensional complex vector space
of convergent holomorphic germs, and let $t>0$.  The set $E_t(F)$ has
exactly $\dim F$ elements, and $F$ admits a basis whose initial exponents
are these distinct elements.

Let $r,m$ be positive integers.  If the monomial space
\[
 \Span_{\C}\{X^\alpha Z^\beta:(\alpha,\beta)\in E_t(F)\}
\]
has an injective order-$m$ jet test at some distinct points
$q_1,\ldots,q_r\in\T$, then representatives of $F$ on a common
neighborhood of the origin have an injective order-$m$ jet test at
distinct points arbitrarily close to the origin.
\end{proposition}

\begin{proof}
The statement is immediate if $F=0$, so write $N=\dim F>0$.

\emph{Choose a basis with distinct initial terms.}
Among the initial exponents of nonzero elements of $F$, choose the
least one, $e_1$.  Every element of $F$ has zero coefficient at every
earlier exponent.  The coefficient at $e_1$ is a nonzero linear
functional on $F$.  Choose $f_1$ on which it equals $1$, and continue
inside its codimension-one kernel.  After $N$ steps this gives a basis
$f_1,\ldots,f_N$ with strictly increasing initial exponents
$e_i=(\alpha_i,\beta_i)$ and initial coefficients $1$.  In a nonzero
linear combination of the basis vectors, the first vector with nonzero
coefficient supplies its uncancelled initial term.  Consequently
$E_t(F)=\{e_1,\ldots,e_N\}$.

\emph{Separate the initial terms by a single positive weight.}
Set
\[
 M=1+\max_{1\le i\le N}(\alpha_i+t\beta_i).
\]
Choose a small number $\eta>0$.  For $\eta$ sufficiently small, all
the weights $\alpha_i+(t+\eta)\beta_i$ remain below $M$.  Every term
of old weight at least $M$ has new weight at least $M$, since its second
exponent is nonnegative, so it cannot become a least term.

Only finitely many exponent pairs have old weight below $M$.  For a
noninitial term of $f_i$ in this finite set, either its old weight is
strictly greater than that of $e_i$, or its old weight is equal and its
second exponent is strictly greater.  In the first case the strict
inequality persists for all sufficiently small $\eta>0$.  In the second
case increasing $t$ by $\eta$ makes the new weight strictly greater.
One common choice of $\eta>0$ therefore makes $e_i$ the unique term of
least weight $(1,t')$ in every $f_i$, where $t'=t+\eta$.  Put
\[
 \lambda_i=\alpha_i+t'\beta_i.
\]

\emph{Dilate the convergent series.}
Choose representatives of the finitely many basis germs on a common
polydisc, and define, for positive real $s$,
\[
 g_{i,s}(X,Z)=s^{-\lambda_i}f_i(sX,s^{t'}Z).
\]
These functions converge locally uniformly on $\C^2$ to
$X^{\alpha_i}Z^{\beta_i}$, in the sense that they are defined on any
fixed compact set for all sufficiently small $s$.  To verify the
convergence, write
$f_i(x,z)=\sum_{\alpha,\beta}c_{i,\alpha,\beta}x^\alpha z^\beta$.
Fix $R>0$ and choose $0<s_0<1$ so that the radii
$(s_0R,s_0^{t'}R)$ lie strictly inside a common polydisc of absolute
Taylor convergence.  Every supported exponent satisfies
$\alpha+t'\beta-\lambda_i\ge0$, with equality only at $e_i$.
For $0<s\le s_0$, the sum of uniform majorants for the rescaled terms
on $|X|,|Z|\le R$ is
\[
 \sum_{\alpha,\beta}
 |c_{i,\alpha,\beta}|s_0^{\alpha+t'\beta-\lambda_i}R^{\alpha+\beta}
 =s_0^{-\lambda_i}
   \sum_{\alpha,\beta}|c_{i,\alpha,\beta}|
     (s_0R)^\alpha(s_0^{t'}R)^\beta
 <\infty.
\]
Each noninitial term tends uniformly to zero.  The summable majorant
therefore proves the asserted uniform convergence.  Applying this
argument on a slightly larger polydisc and using the Cauchy integral
formula gives convergence of all derivatives of any fixed finite
order on the smaller polydisc.

\emph{Transfer the jet test.}
At the given points $q_j$, the limiting monomials have an injective jet
test.  Lemma~\ref{lem:rank-persistence} shows that the functions
$g_{1,s},\ldots,g_{N,s}$ have an injective jet test at the same points
for every sufficiently small positive $s$.  The map
\[
 D_s(X,Z)=(sX,s^{t'}Z)
\]
is invertible for each such $s$.  Its images $p_j=D_s(q_j)$ are distinct,
lie in the common coordinate neighborhood, and approach the origin.
Pullback by $D_s$ preserves orders of vanishing at corresponding
points.  The nonzero scalar factors $s^{-\lambda_i}$ change only the
basis.  Hence the original space $F$ has an injective order-$m$ jet
test at $p_1,\ldots,p_r$.
\end{proof}

The auxiliary weight $t'$ in this proof separates finitely many initial
terms chosen using $t$; those exponents remain unchanged.  Thus a bound
for $E_t(F)$ in a fixed diagram is preserved throughout the argument.
The weights and the positive real dilation parameter need not be
rational, and the construction makes no assertion about their uniformity
over different finite-dimensional spaces.

To apply Proposition~\ref{prop:initial-transfer} to global sections on
$S$, choose local coordinates and a nonvanishing frame for the line
bundle.  Restriction to germs is injective: a section vanishing on a
neighborhood vanishes everywhere by the identity principle on the
connected surface.  A change of frame multiplies each germ by a
holomorphic unit, whose nonzero constant term preserves its initial
exponent.  In particular, if all the initial exponents lie in a bounded
set $\Omega\subset\R_{\ge0}^2$, an injective jet test for the larger
space $\mathcal M(\Omega)$ suffices for the transfer.
